\documentclass[12pt]{article}
\usepackage[utf8]{inputenc}
\usepackage[T1]{fontenc}
\usepackage{hyperref}
\title{Mbudum Language Resources for K-12 Curriculum Development}
\author{Compiled by Claude Sonnet 5.5 (claude-sonnet-5-5) for Andr\'as Kornai}
\date{}
\begin{document}
\maketitle

\section{Introduction}
Mbudum (Mb\textschwa dam, Mbedam; ISO 639-3 xmd; Glottocode mbed1242) is a Chadic (Biu-Mandara, Dabaic) language of Cameroon's Far North region, northeast of Hina in the Mokolo arrondissement. Wikipedia cites about 6,000 speakers in 2001 and says the number is declining. Joshua Project gives 12,000 for the people group. Little material exists, so this report is short.

\section{Primary Educational Resources}
No primers, textbooks or school programmes were found in the search results.

\section{Linguistic and Reference Materials}
\begin{itemize}
\item Wikipedia, Mbudum language: \url{https://en.wikipedia.org/wiki/Mbudum_language}
\item Glottolog reference: Ndokobai, Dadak (2011), a 22-page orthography sketch (described in the listing; the document itself was not seen): \url{https://glottolog.org/resource/reference/id/326050}
\item ASJP word list entry (Mbedam): \url{https://asjp.clld.org/languages/MBEDAM}
\item Glosbe entry for Mbudum (online word lookup; quality unverified): \url{https://glosbe.com/xmd/lib}
\end{itemize}

\section{Religious and Community Resources}
\begin{itemize}
\item Scripture Earth, Mbudum (a New Testament translation is reported): \url{https://www.scriptureearth.org/00eng.php?iso=xmd}
\item Joshua Project, Mbudum language page: \url{https://joshuaproject.net/index.php/languages/xmd}
\end{itemize}

\section{Digital and Audio Resources}
\begin{itemize}
\item Global Recordings Network, Mbudum: \url{https://globalrecordings.net/en/language/xmd}
\end{itemize}

\section{Curriculum Development Context}
The New Testament and the orthography sketch are the main written leads. A school curriculum would need a fuller description of the language and community-led materials. The searches found no existing school use.
\end{document}